\documentclass[11pt]{article}
\usepackage[margin=1in]{geometry}
\usepackage{amsmath,amssymb}
\usepackage{booktabs}
\usepackage{array}
\usepackage{hyperref}
\hypersetup{colorlinks=true,linkcolor=blue,urlcolor=blue}
\emergencystretch=4em

\newcommand{\eps}{\varepsilon}
\newcommand{\epsth}{\eps_\theta}
\newcommand{\epsfr}{\bar\eps}
\newcommand{\xt}{x_t}
\newcommand{\E}{\mathbb{E}}
\newcommand{\norm}[1]{\left\lVert #1 \right\rVert}
\newcommand{\cnull}{\varnothing}

\title{What we train, what the board writes,\\ and the avenues between them}
\author{poe\_repair\_min, thread off \texttt{designing-the-correction-loss}}
\date{2026-09-16}

\begin{document}
\maketitle

\begin{abstract}
Three training objectives for the product-of-experts corrector, written in one notation and
compared as mathematics: the one running in this repository, and the two on the board. They sit on
three independent axes, and separating those axes is most of the work, because two settings that
look like different experiments turn out to be the same experiment written twice. The document
then asks what each objective is really requesting in density terms, what the adapter would have to
have learned for the request to be a reasonable one, and which avenues follow. Measured numbers are
deliberately not repeated here; they live in the two companion documents named at the end.
\end{abstract}

\tableofcontents

\section{Notation, once}
\label{sec:words}

\paragraph{The three branches.}
At denoising step $t$ the network is run three times on the same noisy latent $\xt$: once given the
first prompt, once given the second, once given nothing at all.
\begin{equation}
  a \;=\; \epsth(\xt,t,c_1), \qquad
  b \;=\; \epsth(\xt,t,c_2), \qquad
  u \;=\; \epsth(\xt,t,\cnull).
\end{equation}
Call them the \textbf{cat branch}, the \textbf{dog branch} and the \textbf{empty branch}. A bar
means the adapter is off, so $\epsfr(\cdot,c)$ is the base model and $\bar u$ is its empty branch.

\paragraph{What the adapter adds.}
\begin{equation}
  \phi_1 \;=\; a - \epsfr(\xt,t,c_1), \qquad
  \phi_2 \;=\; b - \epsfr(\xt,t,c_2), \qquad
  m \;=\; u - \bar u .
  \label{eq:changes}
\end{equation}
$\phi_1$ is the \textbf{cat change}, different for every concept. $m$ is the \textbf{shared
change}, the same whichever pair is being trained. These three are the only things training moves.

\paragraph{The composition rule and the guidance dial.}
\begin{equation}
  T \;=\; a + b - u,
  \qquad
  T^{(w)} \;=\; u + w(a-u) + w(b-u) \;=\; w\,T - (w-1)\,u .
  \label{eq:compose}
\end{equation}
Setting $w=1$ returns $T$, so these are one expression at two dial settings
(\texttt{metrics.py:12--17}, combined in \texttt{trainer.py:331--332} at $w = 7.5$). The right-hand form is the one to remember: \textbf{the empty branch is read twice},
once inside the composition and once on its own with coefficient $-(w-1)$.

\paragraph{The target.}
$\eps_J = \epsfr(\xt,t,c_{1\wedge2})$ is what the frozen model predicts when handed the single
joint string \texttt{"a cat and a dog"}. It is read from cache, never recomputed.

\paragraph{Noise space and score space.}
\begin{equation}
  \eps(\xt,t,c) \;=\; -\sigma_t \, \nabla_{\xt} \log p_t(\xt \mid c),
  \label{eq:bridge}
\end{equation}
so adding and subtracting branches is multiplying and dividing densities. The board writes its
losses in score space with a $1/\sigma_t$ on the target; the two spellings differ by $\sigma_t^{-2}$,
which folds into the per-level weight and changes nothing that follows. Everything below is in
noise space.

\paragraph{One event.}
$A$ is the event \emph{this picture contains two separate objects}. So $p(x)$ is
\textbf{all pictures} and $p(x \mid A)$ is \textbf{two-thing pictures}.

\section{The three objectives}

\paragraph{A, what this repository trains.}
\begin{equation}
  \boxed{\;
  \mathcal{L}_A \;=\; \E\;\norm{\;
  \underbrace{T^{(w)}}_{\text{composed, dial at }7.5}
  \;-\;
  \underbrace{\big[\bar u + w(\eps_J - \bar u)\big]}_{\text{joint prompt, dial at }7.5}
  \;}^2
  \;}
  \label{eq:LA}
\end{equation}
Run the adapted model three times, combine, and make the answer match what the frozen model would
have said if simply asked for both animals at once. Both sides carry the dial. One arm adds
$\mu\norm{m}^2$, the \emph{null anchor}, which fines the empty branch for moving
(\texttt{trainer.py:496--503}).

\paragraph{B2, the board with the mono target.}
\begin{equation}
  \boxed{\;
  \mathcal{L}_{B2} \;=\; \E\;\norm{\;\eps_J \;-\; T\;}^2
  \;}
  \label{eq:LB2}
\end{equation}
Same target, no dial anywhere. The board's note beside it: \emph{we assume mono works, PoE must
learn to copy mono.}

\paragraph{B1, the board with the data target.}
\begin{equation}
  \boxed{\;
  \mathcal{L}_{B1} \;=\; \E_{\,x \sim p_{\mathrm{co}},\; z}\;\norm{\;z \;-\; T\;}^2 ,
  \qquad \xt = \text{noised } x
  \;}
  \label{eq:LB1}
\end{equation}
Take a real photograph in which both concepts appear as separate objects, add noise $z$, and ask the
composition to predict the noise you added. $p_{\mathrm{co}}$ is that corpus.

\paragraph{The board's two protection terms.}
They are not the same thing, and the difference matters later.
\begin{align}
  \text{(P2)}\quad &\text{pull a branch toward the \emph{true} score of real single-concept images,}
  \label{eq:P2}\\
  \text{(P3)}\quad &\text{pull a branch toward \emph{the frozen model's own} prediction for that
  concept.}
  \label{eq:P3}
\end{align}
Neither is implemented. The null anchor is (P3) with the empty prompt in place of a concept prompt,
which is why the two are easy to confuse.

\section{Three axes, and only one of them is a real choice between A and B}
\label{sec:axes}

\begin{center}
\small
\begin{tabular}{p{0.16\textwidth} p{0.36\textwidth} p{0.40\textwidth}}
\toprule
axis & settings & who fixes it \\
\midrule
the target & $\eps_J$, or $z$, or something else & A and B2 pick $\eps_J$, B1 picks $z$ \\
the dial & $w = 1$ or $w = 7.5$ & A picks $7.5$, B1 and B2 pick $1$ \\
the freeze pattern & which of $a$, $b$, $u$ carry the adapter & \textbf{nobody: free in all three} \\
\bottomrule
\end{tabular}
\end{center}

The third row is the point of this table. The freeze pattern is not part of B1 and not part of A.
It crosses every target, so anything learned about it under one target transfers to the others.

A and B2 differ on the dial alone. Section~\ref{sec:dial} shows that difference is exactly one term,
and Section~\ref{sec:freeze} shows the freeze pattern can delete it.

\section{The dial: one identity}
\label{sec:dial}

Substitute \eqref{eq:compose} into \eqref{eq:LA}. Everything cancels but two terms:
\begin{equation}
  \boxed{\;
  T^{(w)} - \big[\bar u + w(\eps_J - \bar u)\big]
  \;=\; w\,\big(T - \eps_J\big) \;-\; (w-1)\,m
  \;}
  \label{eq:key}
\end{equation}
so, writing $\Delta = T - \eps_J$ for B2's own error,
\begin{equation}
  \mathcal{L}_A \;=\; w^2\,\mathcal{L}_{B2}
  \;-\; 2w(w-1)\,\langle \Delta, m\rangle
  \;+\; (w-1)^2\,\norm{m}^2 .
  \label{eq:expand}
\end{equation}

Read it in three steps. At the start of training the adapter is zero, so $m = 0$, the last two terms
vanish, and the two objectives are one objective differing by a constant factor $w^2$.

Then the empty branch begins to move. Adding the same vector $v$ to a concept branch and to the
empty branch leaves $a+b-u$ untouched, so $\mathcal{L}_{B2}$ is exactly invariant along that
direction. $\mathcal{L}_A$ is not: its composition shifts by $-(w-1)v$ while its target, built from
the cached $\bar u$, stays put.

So the two objectives disagree about whether that move is progress. \textbf{The dialled loss can
lower itself by shifting the empty branch, while the undialled one records nothing.} Whether the
runs actually took that shortcut is a measurement, and it has been made; see
Section~\ref{sec:numbers}.

\section{What each objective is actually asking for}
\label{sec:asking}

\paragraph{Before training.}
Put the frozen branches through \eqref{eq:bridge}:
\begin{equation}
  \bar T \;=\; -\sigma_t \nabla \log\!\left[\frac{p(x\mid c_1)\,p(x\mid c_2)}{p(x)}\right],
  \qquad
  q_0(x) \;\propto\; \frac{p(x\mid c_1)\,p(x\mid c_2)}{p(x)}
  \;\propto\; p(x)\;p(c_1\mid x)\;p(c_2\mid x),
  \label{eq:q0}
\end{equation}
the last step by Bayes on each conditional. The right-hand form explains the failure. Sampling
$q_0$ means finding something that looks like a real image and that a cat detector \emph{and} a dog
detector both score highly. One fused animal scores well on both across the whole frame; a genuine
two-animal scene gives each detector only part of it. The peak of $q_0$ sits on the fused animal.

\paragraph{One correction of tense, worth making before the paper does not.}
Equation~\eqref{eq:q0} describes $\bar T$, the composition \emph{before} training. The trained
composition is not converging on $q_0$; it is converging on $T = \eps_J$. Stating the fused-animal
argument as a property of the trained adapter is the version a reviewer will catch.

\paragraph{What matching $\eps_J$ demands.}
Bayes gives the exact posterior and the error in one line:
\begin{equation}
  p(x \mid c_1, c_2) \;=\; q_0(x)\cdot R(x),
  \qquad
  R(x) \;=\; \frac{p(c_1,c_2\mid x)}{p(c_1\mid x)\,p(c_2\mid x)}\times(\text{constant in }x).
  \label{eq:R}
\end{equation}
So $a + b - u = \eps_J$ asks for $\nabla_x \log R(x) = 0$: the two concepts must tell you nothing
about each other once you have seen the picture. That is false, and it is false worst at a fused
animal, where both single-concept posteriors are high and the joint one is low, which is exactly
where the training states are.

\paragraph{So say the narrower thing.}
The adapter cannot make two concepts independent, because nothing it does touches
$p(x\mid c_1,c_2)$. It fits three single-prompt functions so that one fixed piece of arithmetic
reproduces a fourth function, on one distribution of states. Two consequences follow, and both are
useful. Matching the target \emph{requires} the branches to move, since frozen branches give the
frozen composition by definition. And the required movement turns out to be per-concept, which is
what makes any of this trainable at all.

\section{The plurality reading: what the adapter would have to have learned}
\label{sec:plurality}

Nothing is written into the prompt. The claim is that $\phi_1$, $\phi_2$ and $m$ are precisely what
turns each branch into its two-thing-picture version:
\begin{equation}
  \text{(H)}\qquad
  a = -\sigma_t \nabla \log p_t(x\mid c_1, A),
  \quad
  b = -\sigma_t \nabla \log p_t(x\mid c_2, A),
  \quad
  u = -\sigma_t \nabla \log p_t(x\mid A).
  \label{eq:H}
\end{equation}

\paragraph{What follows at once.}
\begin{equation}
  T \;=\; a+b-u \;=\; -\sigma_t \nabla \log\!\left[
  \frac{p(x\mid c_1,A)\,p(x\mid c_2,A)}{p(x\mid A)}\right],
  \qquad
  q(x) \;\propto\; \frac{p(x\mid c_1,A)\,p(x\mid c_2,A)}{p(x\mid A)} .
  \label{eq:q}
\end{equation}

\paragraph{Why that is a smaller ask.}
Rerun \eqref{eq:R} with $A$ carried through:
\begin{equation}
  p(x\mid c_1,c_2,A) \;=\; q(x)\cdot R_A(x),
  \qquad
  R_A(x) \;=\; \frac{p(c_1,c_2\mid x,A)}{p(c_1\mid x,A)\,p(c_2\mid x,A)}\times(\text{constant}).
  \label{eq:RA}
\end{equation}
Same shape as \eqref{eq:R}. What changed is the range of $x$ over which the flatness must hold.
$R$ swings hardest at a fused animal; conditioning on $A$ takes fused animals out of the range.
\textbf{The independence is asked only among pictures that already have two objects in them}, which
is a far weaker thing to require and it is required where it is plausible.

\emph{The last sentence is a conjecture, not a derivation.} Nothing has measured $R_A$ against $R$
on this pool. It is measurable.

\paragraph{The catch, and it is a real one.}
The loss only ever sees $a+b-u$. Shift the shared change into the concept changes and the total is
unchanged, so training does not decide how the ``and'' is split between the empty branch and the
concept branches. It decides only the sum. That does not make (H) wrong; it means (H) needs either
a rule that pins the split or a measurement that survives the split being loose.

\section{Which branches to free}
\label{sec:freeze}

Write switches $\mathbf{1}_a,\mathbf{1}_b,\mathbf{1}_u \in \{0,1\}$:
\begin{equation}
  T \;=\; \bar T \;+\; \mathbf{1}_a \phi_1 \;+\; \mathbf{1}_b \phi_2 \;-\; \mathbf{1}_u\, m .
  \label{eq:switches}
\end{equation}

\paragraph{Eight settings, four families.}
$m$ is a function of the picture alone, while $\phi_i$ is a function of the picture \emph{and} the
concept, so $\phi_i - m$ is still a legal per-concept function. Whenever at least one concept
branch is free, the shared change can be absorbed into it and the empty branch's switch changes
nothing:
\begin{equation}
  \{a,b,u\} = \{a,b\}, \qquad
  \{a,u\} = \{a\}, \qquad
  \{b,u\} = \{b\},
  \label{eq:collapse}
\end{equation}
and $\{a\}$ equals $\{b\}$ up to which concept is labelled first. The families nest, and the
least-squares floor runs the other way:
\begin{equation}
  \{u\} \;\subseteq\; \{a\} \;\subseteq\; \{a,b\} \;=\; \{a,b,u\},
  \qquad
  R^{\{u\}} \;\ge\; R^{\{a\}} \;\ge\; R .
\end{equation}

\begin{center}
\small
\begin{tabular}{p{0.19\textwidth} p{0.18\textwidth} c p{0.40\textwidth}}
\toprule
family & free objects & forwards & note \\
\midrule
$\{a,b\}$, $u$ frozen & $\phi_1,\phi_2$ & $2$ & the full family, split pinned \\
$\{a,b,u\}$ & $\phi_1,\phi_2,m$ & $3$ & the same answers, split loose. What runs today \\
$\{a\}$ & $\phi_1$ & $2$ & only one concept may move; no reason to want it \\
$\{u\}$ & $m$ & $1$ & the shared change alone; see Section~\ref{sec:three} \\
\bottomrule
\end{tabular}
\end{center}

So the whole permutation is three experiments, of which two are worth running. Freezing is cheap
because it is not freezing weights: all branches share the same adapted weights, so frozen means
reading a cached frozen tensor, which costs no forward pass.

\paragraph{Two rules that apply to every pattern.}
Whatever is frozen in training must be frozen in sampling. The adapter sits on cross-attention
projections and touches every branch, so training with a frozen empty branch while the sampler runs
all three adapted puts the amplified direction straight back.

And with $m = 0$, equation \eqref{eq:key} loses its second term outright, so
$\mathcal{L}_A = w^2 \mathcal{L}_{B2}$ identically, at every $w$, in value and in gradient.
\textbf{Freezing the empty branch is what makes A and B2 the same objective.} The dial becomes a
constant that folds into the learning rate, and one adapter is correct at every dial setting.

\section{Three ways to place the plurality, and what each gives}
\label{sec:three}

Take (H) seriously and ask where the two-thing conditioning is allowed to live.

\paragraph{Only in the empty branch.}
Freeze the concept branches, let the adapter move only $u$, and suppose it succeeds:
\begin{equation}
  q_c(x) \;\propto\; \frac{p(x\mid c_1)\,p(x\mid c_2)}{p(x\mid A)}
  \;=\; q_0(x)\,\frac{p(x)}{p(x\mid A)}
  \;\propto\; \frac{q_0(x)}{p(A\mid x)} ,
  \label{eq:qc}
\end{equation}
using $p(x\mid A) = p(A\mid x)\,p(x)/p(A)$. You end up \emph{dividing} by the probability that the
picture has two objects in it. A fused animal has small $p(A\mid x)$, so it is boosted; a genuine
two-animal scene has $p(A\mid x)\approx 1$ and is left alone. \textbf{Swapping the denominator alone
does not fix the fused animal, it rewards it.} The direction is wrong, not merely weak.

\paragraph{Only in the concept branches.}
Freeze the empty branch, so $m=0$, and let $a$ and $b$ take their values from \eqref{eq:H}:
\begin{equation}
  q_{\mathrm{freeze}}(x) \;\propto\;
  \frac{p(x\mid c_1,A)\,p(x\mid c_2,A)}{p(x)}
  \;=\; q(x)\,\frac{p(x\mid A)}{p(x)}
  \;\propto\; q(x)\,p(A\mid x) .
  \label{eq:qfreeze}
\end{equation}
The two-thing product, \emph{multiplied} by the probability the picture is a two-thing picture. The
exact mirror of \eqref{eq:qc}, factor upstairs instead of downstairs, leaning toward two-object
pictures instead of away from them.

\paragraph{In all three, as (H) writes it.} You get $q$ itself.

\begin{center}
\begin{tabular}{ll}
\toprule
plurality only in the empty branch & $q_0 \,/\, p(A\mid x)$ \quad rewards fused animals \\
plurality only in the concept branches & $q \cdot p(A\mid x)$ \quad rewards two-thing pictures \\
plurality in all three & $q$ \quad the two-thing product, clean \\
\bottomrule
\end{tabular}
\end{center}

\paragraph{Being fair about the middle row.}
$p(A\mid x)$ is a bias. You are not sampling the two-thing posterior, you are sampling it tilted
once further toward plurality. That is the same shape as guidance, which already tilts hard. It is
a bias in the direction wanted, and it is a bias.

\paragraph{And the middle row is the one that is cheap.}
It is the freeze of Section~\ref{sec:freeze}: two forwards instead of three, the dial removed from
the problem, and the split pinned. The clean bottom row needs the empty branch to move, which
reopens the direction the sampler amplifies, unless the move is supervised rather than left to
drift. That is the avenue in Section~\ref{sec:avenues} called \emph{anchor to plurality}.

\section{Where B1 stands, and what it does not need}

B1 differs from B2 on one axis only, the target. Everything in
Sections~\ref{sec:asking} to~\ref{sec:three} is about the composition and the branches, so it holds
for B1 unchanged.

What B1 needs that the others do not is $p_{\mathrm{co}}$: real photographs containing both concepts
as separate objects. This repository has no real image data at all. The coverage question was
measured against an annotated corpus and the answer was that almost no pair in this pool has such
images, and the few that do are look-alike pairs.

The reason is structural rather than a shortage. \textbf{A pair has co-occurrence photographs
because the two things genuinely turn up together, and that same co-occurrence in the model's
training data is why the model already composes them.} The data is abundant where it is not needed
and absent where it is. Any photographic corpus should behave the same way, though only one has
been checked.

So B1's shape is reachable; B1's corpus is not. The shape survives if the corpus is replaced, which
is the avenue called \emph{a better teacher}.

\section{The avenues}
\label{sec:avenues}

\begin{center}
\small
\begin{tabular}{p{0.17\textwidth} p{0.29\textwidth} p{0.215\textwidth} p{0.235\textwidth}}
\toprule
 & what changes & what it buys & what would kill it \\
\midrule
freeze the empty branch &
$m = 0$, by construction rather than by a fine &
A and B2 become one objective at every dial setting; one fewer forward &
nothing known \\[3pt]
anchor to plurality &
hold $u$ to the frozen prediction for a plurality phrase, not for the empty prompt &
the clean $q$ with the split still pinned &
a plurality \emph{string}'s prediction may be a poor stand-in for $p(x\mid A)$ \\[3pt]
protect the concepts &
add (P3) on states from single-concept runs, not pair runs &
bounded damage to single-concept behaviour &
it opposes (H) unless the concept change depends on where it is measured \\[3pt]
a better teacher &
keep B1's shape, replace $p_{\mathrm{co}}$ with the model's own filtered successes &
a target that is right where the joint prompt is wrong &
the filter cannot tell two animals from one animal twice \\[3pt]
refresh the states &
re-record the cache with the adapter attached, and iterate &
the fit holds where the adapter actually goes &
the off-policy gap turns out to be small \\[3pt]
leave the family &
a read-out seeing both concepts, or a sampler-side correction &
escapes the additive ceiling &
the additive ceiling is not what is binding \\
\bottomrule
\end{tabular}
\end{center}

\subsection*{Freeze the empty branch}
The first move, because it is exact where the current fine is approximate, it costs less than what
runs now, and it removes an axis from the problem. By \eqref{eq:collapse} it gives up no reachable
predictions. By \eqref{eq:qfreeze} it does not cost the plurality reading either; it decides where
the plurality has to live.

\subsection*{Anchor to plurality}
Hold the empty branch to the frozen model's answer for a plurality phrase instead of for the empty
prompt. This pins the split, so the amplified direction stays controlled, \emph{and} puts the base
measure on two-thing pictures rather than waiting for drift to arrive there. It costs one cached
tensor. It also supplies a reference the code currently refuses to guess: the trainer rejects the
anchor together with a non-plain branch style, because the only cached frozen reference is the one
for the empty prompt and anchoring against the wrong reference is worse than not anchoring
(\texttt{train\_pooled.py:315--321}).

The honest caveat: a frozen prediction for the string \texttt{"two animals"} is a prediction for a
string, not the score of $p(x \mid A)$, and how close those are is unchecked.

\subsection*{Protect the concepts}
Three candidates, one recommendation. The null anchor already in the code is not single-concept
protection: it constrains the one branch that carries no concept, and its job is pinning the split.
(P2) pulls a branch toward the truth rather than toward where it started, which is fine-tuning, not
preservation, and it needs real single-concept images. (P3) needs no images, only states, and the
choice of states is the whole design.

Applied on the pair states, (P3) opposes the fit directly, because the fit requires those branches
to move there. Push it hard and all three branches pin to frozen and the correction goes to zero.
So apply it on states from single-concept runs instead: correct the model where pairs go, preserve
it where a single concept goes.

\paragraph{One condition, and it is not optional.}
Under (H), this term pulls the concept branches back from the very thing they are supposed to
become. The two goals survive together only if $\phi_1$ is large where pairs go and small where a
single concept goes. Settle that before the term goes near a training run, not after.

\subsection*{A better teacher}
Keep the shape of \eqref{eq:LB1} and replace $p_{\mathrm{co}}$ with the model's own samples that
pass a composition filter. The reason to want it is that the current teacher is known to be
unreliable: the scope that trains the pooled adapter already drops training cells whose joint target
shows the wrong animals, and names a better teacher than the joint prompt as its own follow-on. A
filtered sample set is right by construction exactly where the joint prompt is wrong.

The risk is the filter. The validated scorer counts instances and cannot tell a cat beside a dog
from two dogs, so the corpus needs a second head reading which two animals are present, or the
objective trains toward a density of duplicated single concepts.

\subsection*{Refresh the states}
Every cached state came from a plain product-of-experts run with no adapter. With the adapter on,
the trajectory leaves those states on the first step it acts, and errors compound along a rollout.
The answer, if it is needed, is to re-record with the adapter attached and iterate. Whether it is
needed is a cheap read, not an assumption: take one adapted step from a cached state and compare the
loss where you land against the loss at the cached successor.

\subsection*{Leave the family}
Every objective above has the shape $T = m + \phi_1 + \phi_2$, one shared piece plus one per
concept, forced by each branch reading one prompt under a fixed combination rule. Escaping it means
a read-out that sees both concept embeddings at once, which costs the property that the adapter
never sees a joint prompt, or a correction on the sampler side. Neither is earned while the
per-concept fit is still doing well.

\section{What would decide between them}

Three reads, none of which needs training, all against adapters already on disk.

\paragraph{Does the cat change point where the word points?}
Compare $\phi_1$, what the adapter added to the cat branch, against what the word adds with no
adapter anywhere: the frozen model's answer for \texttt{"a cat and"} minus its answer for
\texttt{"a cat"}. Report the cosine. The control is the same quantity against another concept's word
direction and against an unrelated word, so a meaningless number can be recognised. Beating the
control supports (H); sitting inside it says the adapter learned something real that is not the
``and''.

Run at two sets of states, pair runs and single-concept runs, this is also what decides whether
\emph{protect the concepts} is safe.

\paragraph{Is the change the same thing for every concept?}
If ``and'' is one learned thing, the per-concept changes share a component. Split $\{\phi_i\}$ into
their mean plus a residual and report the mean's share of the energy. This costs nothing: the
existing step-$0$ probe already solves for $\phi$ and then discards it.

\paragraph{Did the prompt version and the plain version land in the same place?}
Writing plurality into the prompt has already been run, twice, at full length: once with
\texttt{"a cat and"} against a null of \texttt{"and"}, once with \texttt{", two animals"} against a
null of \texttt{"two animals"}. By \eqref{eq:collapse} the plain setup can express whatever those
express. Whether training actually found the same place is a fact about optimisation, and it is
unmeasured. Compare the adapted plain branch for \texttt{"a cat"} against the adapted connective
branch for \texttt{"a cat and"} at the same state.

\section{Where the numbers are}
\label{sec:numbers}

Deliberately not repeated here. Two companion documents in this folder carry them with their
sources:

\begin{itemize}
  \item \texttt{guided-and-unguided-objectives.tex}: how large the shared change actually grew on
    each trained arm, what that does to \eqref{eq:expand}, and the bracket it puts on
    $\mathcal{L}_{B2}$ for the plain and anchored runs. Also the held-out per-concept fit, the
    corpus coverage count, and the probe for the off-policy gap with its bar.
  \item \texttt{correction-loss-variants.tex}: the additive skeleton and why the branch prompts can
    be renamed without changing what is reachable, the step-$0$ floor, and the measurements that
    bear on the shared change acting alone.
\end{itemize}

\end{document}
